% -----------------------------------------------------------------------------
% Inhaltsteil zum direkten Einbinden in die Masterarbeit
% Benoetigte Pakete:
% booktabs, tabularx, longtable, array, ragged2e, xcolor, tcolorbox,
% enumitem, hyperref, microtype
% -----------------------------------------------------------------------------

\section{Systemfunktionen und Feature-Abdeckung}
\label{sec:systemfunktionen_featureabdeckung}

\subsection{Zweck der funktionsorientierten Sicht}

Die bisherige Darstellung nach Logical Components eignet sich insbesondere zur
Beschreibung der inneren Systemarchitektur: Sie zeigt, welche logischen
Verantwortungsbereiche vorgesehen sind und wie diese miteinander interagieren.
Für die Beantwortung der Frage, \emph{was das System fachlich leistet}, ist eine
funktionsorientierte Sicht jedoch verständlicher.

Die nachfolgende Feature-Übersicht wird deshalb anhand der zwölf
Systemfunktionen strukturiert. Jede Funktion beschreibt einen abgegrenzten
fachlichen Beitrag des Turing Generators. Die Logical Components bleiben
weiterhin relevant, dienen hier jedoch nur als Realisierungsstruktur und nicht
als primäre Gliederung der Features.

\begin{tcolorbox}[
  enhanced,
  breakable,
  colback=blue!3,
  colframe=TuringBlue,
  title={Lesart der folgenden Darstellung},
  fonttitle=\bfseries,
  boxrule=0.8pt,
  arc=2mm
]
Die Systemfunktionen beschreiben den vollständigen angestrebten Systemumfang.
Der Implementierungsstatus zeigt dagegen, welcher Anteil davon in der
aktuellen Phase-F/P-Baseline bereits umgesetzt und verifiziert ist. Eine
Funktion kann daher bereits einen demonstrierbaren Kern besitzen und trotzdem
als \emph{teilweise implementiert} bewertet werden, wenn ihr vollständiger
Zielumfang noch nicht erreicht ist.
\end{tcolorbox}

\subsection{Was macht das System?}

Auf einer hohen Abstraktionsebene verarbeitet der Turing Generator
heterogene technische Bestandsinformationen in vier zusammenhängenden
Funktionsbereichen.

\subsubsection*{1. Kontext schaffen und Eingangsinformationen bewerten}

Zunächst verwaltet das System Projekte und die ihnen zugeordneten Quellen.
Anschließend werden die verfügbaren Informationen extrahiert, auf ihre
Verarbeitbarkeit geprüft, semantisch eingeordnet und hinsichtlich ihrer
Abdeckung bewertet.

\begin{itemize}[leftmargin=7mm,itemsep=1mm]
  \item \textbf{SF\_001} verwaltet Projekte, Quellen und deren Kontext.
  \item \textbf{SF\_002} bewertet und bereitet Quellinformationen auf.
  \item \textbf{SF\_003} sorgt für eine konsistente technische Bedeutung.
  \item \textbf{SF\_004} bewertet, welche Modellbereiche durch Evidenz gestützt sind.
\end{itemize}

\subsubsection*{2. Verarbeitung steuern und menschliche Entscheidungen einbinden}

Die Verarbeitung wird nicht als starre lineare Kette ausgeführt. Das System
konfiguriert Processing Runs, erzeugt bei Bedarf alternative fachliche
Interpretationen und bindet den Menschen an den entscheidenden Stellen als
Review- und Freigabeinstanz ein.

\begin{itemize}[leftmargin=7mm,itemsep=1mm]
  \item \textbf{SF\_005} konfiguriert und orchestriert Processing Runs.
  \item \textbf{SF\_006} erzeugt und vergleicht fachliche Kandidaten.
  \item \textbf{SF\_007} unterstützt Review, Konfliktlösung und Freigabe.
\end{itemize}

\subsubsection*{3. Architektur ableiten, validieren und als SysML v2 bereitstellen}

Aus geprüften Informationen sollen strukturierte Architekturinformationen
abgeleitet werden. Diese werden gegen Struktur-, Konsistenz- und
Integrationskriterien geprüft und schließlich als textuelle SysML-v2-Artefakte
bereitgestellt.

\begin{itemize}[leftmargin=7mm,itemsep=1mm]
  \item \textbf{SF\_008} leitet Architekturinformation ab und strukturiert sie.
  \item \textbf{SF\_009} validiert und integriert Architekturinhalt.
  \item \textbf{SF\_010} erzeugt und prüft SysML-v2-Architekturartefakte.
\end{itemize}

\subsubsection*{4. Evidenz, Traceability und Transparenz sicherstellen}

Zwei Funktionen wirken über den gesamten Verarbeitungsweg hinweg. Sie sorgen
dafür, dass Ergebnisse nicht nur erzeugt, sondern auch nachvollzogen,
rekonstruiert und für den Nutzer verständlich dargestellt werden können.

\begin{itemize}[leftmargin=7mm,itemsep=1mm]
  \item \textbf{SF\_011} erhält Evidenz, Provenienz und Traceability.
  \item \textbf{SF\_012} stellt Status, Review-Bedarf und Historie dar.
\end{itemize}

\begin{tcolorbox}[
  enhanced,
  breakable,
  colback=gray!4,
  colframe=TuringSlate,
  title={Zusammenspiel der Funktionen},
  fonttitle=\bfseries,
  arc=2mm
]
Die Funktionen bilden kein reines Fließband. Semantische Konflikte,
unzureichende Evidenz, Review-Entscheidungen oder Validierungsbefunde können
Rücksprünge in frühere Funktionen auslösen. SF\_005 koordiniert diese
Verarbeitung, während SF\_011 und SF\_012 den gesamten Ablauf mit Evidenz,
Traceability und Transparenz unterstützen.
\end{tcolorbox}
\pagebreak
\subsection{Kompakte Funktionsübersicht}

\renewcommand{\arraystretch}{1.18}
\small
\begin{longtable}{
  >{\RaggedRight\arraybackslash}p{1.45cm}
  >{\RaggedRight\arraybackslash}p{4.1cm}
  >{\RaggedRight\arraybackslash}p{6.3cm}
  >{\RaggedRight\arraybackslash}p{2.8cm}
}
\caption{Funktionsorientierte Übersicht des Turing Generators}
\label{tab:function_summary}\\
\toprule
\textbf{ID} & \textbf{Systemfunktion} & \textbf{Verständliche Kurzbeschreibung} & \textbf{Aktueller Status} \\
\midrule
\endfirsthead
\toprule
\textbf{ID} & \textbf{Systemfunktion} & \textbf{Verständliche Kurzbeschreibung} & \textbf{Aktueller Status} \\
\midrule
\endhead
\midrule
\multicolumn{4}{r}{\footnotesize Fortsetzung auf der nächsten Seite}\\
\endfoot
\bottomrule
\endlastfoot

SF\_001 & Manage Engineering Project and Source Context &
Verwaltet getrennte Engineering-Projekte und die zugehörigen Quellen, sodass
jede Verarbeitung im richtigen fachlichen Kontext erfolgt. &
\statuspartial{Teilweise implementiert} \\

SF\_002 & Assess and Prepare Source Information &
Liest unterstützte Quellen aus, bewertet ihre Qualität und bereitet die
verwendbaren Informationen für die weitere Verarbeitung vor. &
\statuspartial{Teilweise implementiert} \\

SF\_003 & Apply Semantic Governance &
Gleicht Begriffe und Bedeutungen mit Vokabularen und Ontologien ab und macht
semantische Konflikte sichtbar. &
\statuspartial{Teilweise implementiert} \\

SF\_004 & Assess Coverage and Evidence Sufficiency &
Zeigt, welche Zielbereiche durch vorhandene Informationen gestützt sind und
wo belastbare Evidenz fehlt. &
\statuspartial{Teilweise implementiert} \\

SF\_005 & Configure and Orchestrate Processing Runs &
Konfiguriert Verarbeitungsläufe, steuert ihre Schritte und erzwingt notwendige
Qualitäts- und Review-Gates. &
\statuspartial{Teilweise implementiert} \\

SF\_006 & Generate and Compare Engineering Candidates &
Erzeugt plausible fachliche Interpretationen, vergleicht Alternativen und
bereitet Unterschiede für den Review auf. &
\statusdone{Implementiert und verifiziert} \\

SF\_007 & Support Human Review and Approval &
Ermöglicht dem Menschen, Ergebnisse zu prüfen, anzupassen, auszuwählen,
abzulehnen und nachvollziehbar freizugeben. &
\statuspartial{Teilweise implementiert} \\

SF\_008 & Derive and Structure Architecture Information &
Überführt geprüfte Informationen in strukturierte, modulare
Architekturkandidaten. &
\statusopen{Architektur definiert} \\

SF\_009 & Validate and Integrate Architecture Content &
Prüft Architekturinhalt auf Konsistenz, Randbedingungen und
Integrationsfähigkeit. &
\statusopen{Architektur definiert} \\

SF\_010 & Generate and Validate SysML v2 Architecture Artifacts &
Erzeugt traceable Architecture-as-Code-Artefakte in SysML-v2-Notation und
prüft deren textuelle Konformität. &
\statusopen{Architektur definiert} \\

SF\_011 & Maintain Processing Evidence and Traceability &
Hält Quellen, Verarbeitungsschritte, Entscheidungen, Ergebnisse und Artefakte
durchgängig nachvollziehbar. &
\statusdone{Implementiert und verifiziert} \\

SF\_012 & Present Processing Status and History &
Zeigt dem Nutzer den aktuellen Zustand, offene Review-Bedarfe, Auffälligkeiten
und die bisherige Verarbeitungshistorie. &
\statusdone{Implementiert und verifiziert} \\

\end{longtable}
\normalsize

\subsection{Feature- und Stakeholder-Requirement-Abdeckung nach Systemfunktion}

Die nachfolgenden Funktionsprofile verbinden vier Perspektiven:

\begin{enumerate}[leftmargin=8mm,itemsep=1mm]
  \item die fachliche Bedeutung der Systemfunktion,
  \item die durch sie abgedeckten Stakeholder Requirements,
  \item den aktuellen Implementierungsstand,
  \item den noch offenen Zielumfang.
\end{enumerate}

Die Zuordnung der Stakeholder Requirements wurde transitiv über die
System Requirements ermittelt. Die System Requirements bleiben in dieser
Darstellung bewusst ausgeblendet, da sie für die fachliche Übersicht zu
detailliert wären.

\newcommand{\featureprofile}[7]{%
\begin{tcolorbox}[
  enhanced,
  breakable,
  colback=white,
  colframe=TuringBlue,
  title={#1 -- #2},
  fonttitle=\bfseries,
  boxrule=0.8pt,
  arc=2mm,
  before skip=4mm,
  after skip=4mm
]
\textbf{Fachlicher Zweck}\par
#3

\medskip
\textbf{Abgedeckte Stakeholder Requirements}\par
#4

\medskip
\begin{tabularx}{\linewidth}{>{\bfseries\RaggedRight\arraybackslash}p{3.2cm}X}
Status & #5\\
Bereits vorhanden & #6\\
Verbleibender Zielumfang & #7\\
\end{tabularx}
\end{tcolorbox}
}

\featureprofile
{SF\_001}
{Manage Engineering Project and Source Context}
{Die Funktion stellt sicher, dass jede Verarbeitung einem eindeutigen
Engineering-Projekt und den richtigen Quellen zugeordnet ist. Sie verhindert
die Vermischung verschiedener Projekte und schafft den Kontext für alle
weiteren Verarbeitungsschritte.}
{\textbf{SR\_005} Source Relevance Classification;
\textbf{SR\_019} Source Traceability;
\textbf{SR\_029} Multiple Project Support;
\textbf{SR\_030} Project Information Isolation;
\textbf{SR\_038} Multiple Source Support;
\textbf{SR\_039} Processing History Traceability.}
{\statuspartial{Teilweise implementiert}; der zentrale Projekt- und
Quellenkontext ist demonstrierbar.}
{Persistente Projekt-Workspaces und Manifeste, sechsstellige Projektidentität,
Projektanlage und Wiederöffnung, Projektisolation, projektlokale
Quellenregistrierung, Source IDs, Rollen, Hashes, Dublettenprüfung,
Projektauswahl und Source Upload.}
{Gemeinsame Verarbeitung mehrerer Quellen innerhalb eines Engineering Runs,
Bearbeiten bestehender Projekte sowie kontrolliertes Löschen von Projekten und
zugeordneten Daten.}
\pagebreak
\featureprofile
{SF\_002}
{Assess and Prepare Source Information}
{Die Funktion macht heterogene Quellen technisch und fachlich nutzbar. Sie
extrahiert Informationen, bewertet deren Qualität und identifiziert Lücken,
Widersprüche oder Randbedingungsverletzungen vor der weiteren Nutzung.}
{\textbf{SR\_001} Input Data Suitability Assessment;
\textbf{SR\_002} Input Data Completeness Check;
\textbf{SR\_003} Input Data Consistency Check;
\textbf{SR\_004} Source Reliability Classification;
\textbf{SR\_005} Source Relevance Classification;
\textbf{SR\_014} Explanation of Generated Engineering Results;
\textbf{SR\_019} Source Traceability;
\textbf{SR\_023} Input Constraint Violation Detection;
\textbf{SR\_026} Legacy Data Completion;
\textbf{SR\_034} Insufficient Evidence Identification;
\textbf{SR\_037} Non-Textual Engineering Information;
\textbf{SR\_038} Multiple Source Support.}
{\statuspartial{Teilweise implementiert}; die textuelle Ingestion und
Bewertung sind demonstrierbar.}
{Deterministische Source Projections für Text, Markdown, JSON, CSV, TSV und
PDF-Textschichten; Phase-F-Assessment mit Gaps, Ambiguities, Risks und
Review Questions; projekt- und rungebundene Verarbeitungsvorbereitung.}
{Kontrollierte Datenvervollständigung, vollständige
quellenübergreifende Konsistenzbewertung sowie OCR,
Zeichnungsinterpretation und allgemeine multimodale Extraktion.}

\featureprofile
{SF\_003}
{Apply Semantic Governance}
{Die Funktion sorgt dafür, dass gleiche oder ähnliche Begriffe konsistent
interpretiert werden. Sie bindet Vokabulare und Ontologien ein und macht
semantische Konflikte für nachfolgende Entscheidungen sichtbar.}
{\textbf{SR\_011} Ambiguity Detection;
\textbf{SR\_014} Explanation of Generated Engineering Results;
\textbf{SR\_016} Technical Vocabulary Alignment;
\textbf{SR\_017} Semantic Equivalence Handling;
\textbf{SR\_018} Ontology Conflict Detection.}
{\statuspartial{Teilweise implementiert}; die semantische Infrastruktur ist
vorhanden, aber noch nicht über den gesamten Zielprozess geschlossen.}
{Versionierte Ontologie-Snapshots, Ontology Registry, Turing Core Vocabulary,
Glossary Candidates, Terminology Decisions, Terminology Mappings und
semantischer Multi-Agenten-Konsens mit sichtbarer Varianz.}
{Vollständige Ontologie-Konfliktlösung und iterative semantische Rückkopplung
aus späteren Architektur- und Modellvalidierungen.}
\pagebreak
\featureprofile
{SF\_004}
{Assess Coverage and Evidence Sufficiency}
{Die Funktion beantwortet, welche Teile des angestrebten Modells durch
vorhandene Informationen gestützt sind. Sie verhindert, dass aus unzureichender
Evidenz unzulässige technische Schlussfolgerungen entstehen.}
{\textbf{SR\_014} Explanation of Generated Engineering Results;
\textbf{SR\_033} Engineering Information Coverage;
\textbf{SR\_034} Insufficient Evidence Identification;
\textbf{SR\_036} Approved Information Boundary.}
{\statuspartial{Teilweise implementiert}; Preliminary Coverage ist
implementiert, Approved Generation Readiness noch nicht.}
{Deterministische Preliminary-Coverage-Bewertung, Support Profile,
Unterscheidung von covered, uncovered und attention required sowie
Darstellung der Coverage-Evidence im Dashboard.}
{Verbindliche Evidence-Sufficiency-Entscheidung für modellbildende Schritte
und die spätere Approved Generation Readiness.}

\featureprofile
{SF\_005}
{Configure and Orchestrate Processing Runs}
{Die Funktion bildet das operative Rückgrat des Systems. Sie bindet Kontext,
Rezepte und Agentenprofile, startet Verarbeitungsschritte und kontrolliert
Zustände, Gates, Retry, Recovery und Nachfolgeläufe.}
{\textbf{SR\_001} Input Data Suitability Assessment;
\textbf{SR\_002} Input Data Completeness Check;
\textbf{SR\_003} Input Data Consistency Check;
\textbf{SR\_004} Source Reliability Classification;
\textbf{SR\_005} Source Relevance Classification;
\textbf{SR\_012} Expert Review of Generated Interpretations;
\textbf{SR\_030} Project Information Isolation;
\textbf{SR\_031} Processing State Transparency;
\textbf{SR\_038} Multiple Source Support;
\textbf{SR\_039} Processing History Traceability.}
{\statuspartial{Teilweise implementiert}; Orchestration bis
\texttt{awaiting\_review} ist demonstrierbar.}
{Immutable Run-, Event- und Decision-Manifeste, deterministische
Zustandsrekonstruktion, Stage-Steuerung, Attempts, Retry, Supersession,
Invalidation, Recovery-Diagnostik und projektgebundene P9-Orchestration.}
{Orchestration der späteren Architektur- und SysML-v2-Stufen,
bedienbare Retry-/Successor-Workflows und operative Performancebewertung
vollständiger Live-Agententeams.}

\featureprofile
{SF\_006}
{Generate and Compare Engineering Candidates}
{Die Funktion erzeugt bei mehrdeutiger oder interpretationsbedürftiger
Information mehrere fachliche Kandidaten. Sie stellt Gemeinsamkeiten,
Unterschiede, Annahmen und Evidenz so bereit, dass ein belastbarer Review
möglich ist.}
{\textbf{SR\_011} Ambiguity Detection;
\textbf{SR\_013} Human Intervention for Conflicts;
\textbf{SR\_014} Explanation of Generated Engineering Results.}
{\statusdone{Implementiert und verifiziert}.}
{Modulare Agenten- und Teamausführung, verschiedene Personas und Runs,
deterministischer Konsens, Varianz- und Disagreement-Evidence, Agent Outputs,
Engineering Candidates und Review Reports.}
{Weitergehende domänenspezifische Ranking- und Auflösungsstrategien.
Architektur- beziehungsweise Modellkandidaten gehören bewusst zu SF\_008.}

\featureprofile
{SF\_007}
{Support Human Review and Approval}
{Die Funktion hält den Menschen als fachliche Autorität im Prozess. Ergebnisse
können geprüft, ausgewählt, kombiniert, verändert oder abgelehnt werden;
Entscheidungen bleiben dabei an den tatsächlich geprüften Inhalt gebunden.}
{\textbf{SR\_012} Expert Review of Generated Interpretations;
\textbf{SR\_013} Human Intervention for Conflicts;
\textbf{SR\_034} Insufficient Evidence Identification;
\textbf{SR\_035} Human Review Decision Association;
\textbf{SR\_036} Approved Information Boundary;
\textbf{SR\_039} Processing History Traceability.}
{\statuspartial{Teilweise implementiert}; Review Requests und Decision Binding
existieren, die fachliche Promotion ist noch offen.}
{Immutable Human Review Decisions mit Target- und Fingerprint-Bindung,
Publication Gates, \texttt{review\_requested}-Events sowie prominente
Navigation zu reviewpflichtigen Ergebnissen.}
{Approved Input Promotion und spätere modellbezogene Review- und
Freigabeschritte bis zur autoritativen Architektur.}
\pagebreak
\featureprofile
{SF\_008}
{Derive and Structure Architecture Information}
{Die Funktion bildet die Brücke von geprüfter Engineering-Information zu
strukturierten Architekturmodellen. Sie ordnet Inhalte dem Zielmodell zu und
formt daraus modulare, wiederverwendbare Architektureinheiten.}
{\textbf{SR\_006} Architecture Information Derivation;
\textbf{SR\_007} Structural Pattern Conformance;
\textbf{SR\_008} Model Reusability;
\textbf{SR\_009} Modular Model Growth;
\textbf{SR\_010} Non Destructive Submodel Integration;
\textbf{SR\_022} Architectural Constraint Conformance;
\textbf{SR\_024} Interface Definition Consistency;
\textbf{SR\_025} Reusable Architectural Pattern Application;
\textbf{SR\_036} Approved Information Boundary.}
{\statusopen{Architektur definiert}; vorbereitende Grundlagen sind vorhanden,
die zentrale Runtime-Funktion jedoch noch nicht.}
{Versioniertes RFLP-Framework, stabile Mapping Targets,
Framework-Assignment-Candidates und Preliminary-Support-Bewertung.}
{Architecture Candidate Layer, modulare Model Units, Constraint- und
Pattern-Anwendung, Interface-Antizipation und eigentliche
Architekturableitung.}

\featureprofile
{SF\_009}
{Validate and Integrate Architecture Content}
{Die Funktion prüft, ob abgeleiteter Architekturinhalt strukturell konsistent
ist und ohne unerwünschte Veränderungen in einen größeren Modellkontext
integriert werden kann.}
{\textbf{SR\_007} Structural Pattern Conformance;
\textbf{SR\_010} Non Destructive Submodel Integration;
\textbf{SR\_015} Model Content Validation;
\textbf{SR\_020} Integrated Model Consistency;
\textbf{SR\_021} Larger Context Compatibility;
\textbf{SR\_022} Architectural Constraint Conformance;
\textbf{SR\_024} Interface Definition Consistency;
\textbf{SR\_025} Reusable Architectural Pattern Application.}
{\statusopen{Architektur definiert}; noch keine modellbezogene
Runtime-Implementierung.}
{Normative Anforderungen, Logical Architecture und wiederverwendbare
Artefakt- und Validierungsmechanismen als Grundlage.}
{Konsistenzprüfung erzeugter Architektur, Larger-Context Compatibility,
non-destruktive Integration sowie Constraint-, Pattern- und
Interface-Konfliktvalidierung.}

\featureprofile
{SF\_010}
{Generate and Validate SysML v2 Architecture Artifacts}
{Die Funktion stellt das Ergebnis als maschinenverarbeitbare
Architecture-as-Code-Artefakte bereit. Sie erzeugt SysML-v2-Text und prüft,
ob Notation und Artefaktstruktur dem gewählten Zielprofil entsprechen.}
{\textbf{SR\_007} Structural Pattern Conformance;
\textbf{SR\_014} Explanation of Generated Engineering Results;
\textbf{SR\_015} Model Content Validation;
\textbf{SR\_019} Source Traceability;
\textbf{SR\_027} Architecture as Code Provision;
\textbf{SR\_028} SysML v2 Target Representation;
\textbf{SR\_036} Approved Information Boundary.}
{\statusopen{Architektur definiert}; die Runtime-Erzeugung ist noch offen.}
{Versionierte Target-Notation- und Target-Structure-Annahmen sowie manuelle
CATIA-/SysML-v2-Modellierung als Design Evidence.}
{Automatisierte SysML-v2-Codegenerierung, Notationsprüfung, Export,
Validierung in der Zielumgebung und CATIA-Synchronisation.}

\featureprofile
{SF\_011}
{Maintain Processing Evidence and Traceability}
{Die Funktion sorgt dafür, dass jederzeit nachvollzogen werden kann, aus
welcher Quelle ein Ergebnis stammt, wie es verarbeitet wurde und welche
Entscheidungen oder Artefakte daraus entstanden sind.}
{\textbf{SR\_017} Semantic Equivalence Handling;
\textbf{SR\_019} Source Traceability;
\textbf{SR\_026} Legacy Data Completion;
\textbf{SR\_030} Project Information Isolation;
\textbf{SR\_031} Processing State Transparency;
\textbf{SR\_032} Review Need Transparency;
\textbf{SR\_033} Engineering Information Coverage;
\textbf{SR\_035} Human Review Decision Association;
\textbf{SR\_037} Non-Textual Engineering Information;
\textbf{SR\_039} Processing History Traceability.}
{\statusdone{Implementiert und verifiziert}.}
{Immutable Source-, Run-, Event-, Decision- und Artifact Records,
SHA-256-Integrität, Provenienz, Event Chains, Source-/Project-Aggregation,
Evidence Navigation und Rekonstruktion der Verarbeitungshistorie.}
{Erweiterung der Traceability über Approved Input und generierte
Modellartefakte bis zur finalen CATIA-Synchronisation.}
\pagebreak
\featureprofile
{SF\_012}
{Present Processing Status and History}
{Die Funktion übersetzt den technischen Systemzustand in eine für Nutzer
verständliche Übersicht. Sie zeigt Fortschritt, offene Reviews,
Attention Conditions, Evidenz und Historie.}
{\textbf{SR\_012} Expert Review of Generated Interpretations;
\textbf{SR\_014} Explanation of Generated Engineering Results;
\textbf{SR\_031} Processing State Transparency;
\textbf{SR\_032} Review Need Transparency;
\textbf{SR\_039} Processing History Traceability.}
{\statusdone{Implementiert und verifiziert}.}
{Project Dashboard mit Overview, Sources \& Processing, Coverage \& Support,
Attention \& Review und Traceability Views; Status, Stage, Review-Bedarf,
Historie und Evidence Preview.}
{Verfeinerte Aktionen für nicht-primäre Evidenz sowie Bedienoberflächen für
Retry, Recovery und Successor Runs.}

\subsection{Gesamtbewertung der Funktionsabdeckung}

Alle 39 Stakeholder Requirements werden durch mindestens eine der zwölf
Systemfunktionen abgedeckt. Diese vollständige modellseitige Abdeckung bedeutet
jedoch nicht, dass alle Funktionen bereits vollständig implementiert sind.

\begin{table}[htbp]
\centering
\caption{Zusammenfassung des aktuellen Funktionsstatus}
\label{tab:function_status_summary}
\begin{tabular}{lcc}
\toprule
\textbf{Status} & \textbf{Anzahl Funktionen} & \textbf{Systemfunktionen} \\
\midrule
Implementiert und verifiziert & 3 & SF\_006, SF\_011, SF\_012 \\
Teilweise implementiert & 6 & SF\_001--SF\_005, SF\_007 \\
Architektur definiert & 3 & SF\_008--SF\_010 \\
\bottomrule
\end{tabular}
\end{table}

Der aktuelle Prototyp besitzt damit eine belastbare Basis für
Projekt- und Quellenverwaltung, textuelle Informationsverarbeitung,
agentische Kandidatenerzeugung, Processing Runs, Coverage, Evidenz,
Traceability und Statusdarstellung. Der noch offene Zielabschnitt beginnt beim
Übergang von geprüfter Engineering-Information zu autoritativen
Architekturkandidaten und reicht bis zur validierten SysML-v2-Erzeugung.

\begin{tcolorbox}[
  enhanced,
  breakable,
  colback=gray!4,
  colframe=TuringSlate,
  title={Hinweis für die finale Fassung der Masterarbeit},
  fonttitle=\bfseries,
  arc=2mm
]
Nach Abschluss des Systems sollte insbesondere der Status von SF\_007 bis
SF\_010 aktualisiert werden. Zusätzlich sollten die finalen Evaluationsergebnisse
zur Qualität der Architekturableitung, zur Modellvalidität, zur Modularität und
zur syntaktischen sowie semantischen Konformität der erzeugten
SysML-v2-Artefakte ergänzt werden.
\end{tcolorbox}
